% ----- CHAUPAI 24 -----

\newpage

\subsection*{\deva{चतुर्विंशतितम चौपाई} \(Twenty-Fourth Chaupai\)}
\addcontentsline{toc}{subsection}{\deva{चतुर्विंशतितम चौपाई} \(Twenty-Fourth Chaupai\) — \deva{भूत पिसाच...}}

\begin{minipage}[t]{0.55\textwidth}
\vspace{0pt}

{\large \linespread{1.5}\selectfont
\deva{भूत पिसाच निकट नहिं आवै।}\\
\deva{महाबीर जब नाम सुनावै॥}
\par}

\vspace{0.5em}

{\large \linespread{1.5}\selectfont
\telu{భూత పిసాచ నికట నహిఁ ఆవై।}\\
\telu{మహాబీర జబ నామ సునావై॥}
\par}

\vspace{0.5em}

{\large \linespread{1.3}\selectfont
\textit{bhūta pisāca nikaṭa nahĩ āvai |}\\
\textit{mahābīra jaba nāma sunāvai ||}
\par}
\end{minipage}%
\hfill
\begin{minipage}[t]{0.42\textwidth}
\vspace{0pt}
\begin{tcolorbox}[anchorbox]
\textbf{Mental Armor:} The "ghosts" (\textit{Bhuta Pisacha}) represent performance anxiety and self-doubt. In the *Valmiki Ramayana* (*Sundara Kanda, 50:14-16*), when standing captive alone in Ravana's hostile court, Hanuman maintains complete composure. He objectively observes the grandeur of his enemies without being intimidated or losing focus. Keeping strong role models in mind acts as mental armor against external pressures.
\vspace{0.5em}\newline Standing alone in a hostile court, Hanuman's objective composure acted as mental armor, keeping him focused and unintimidated.\newline\null\hfill\href{https://www.valmiki.iitk.ac.in/sloka?field_kanda_tid=5&language=dv&field_sarga_value=50}{\textit{Sundara Kanda, 50:14-16}}
\end{tcolorbox}
\end{minipage}

\vspace{0.5em}
\subsubsection*{\textbf{Visual Rhythm Map:}}
\begin{table}[h!]
\centering
\renewcommand{\arraystretch}{1.2}
\setlength{\tabcolsep}{0pt}
\begin{tabularx}{\textwidth}{|*{16}{>{\centering\arraybackslash\hsize=1.0\hsize}X|}}
\hline
\multicolumn{3}{|>{\centering\arraybackslash\hsize=3.0\hsize}X|}{\textbf{\guru{भू}\laghu{त}} \par (3)} &
\multicolumn{4}{>{\centering\arraybackslash\hsize=4.0\hsize}X|}{\textbf{\laghu{पि}\guru{सा}\laghu{च}} \par (4)} &
\multicolumn{3}{>{\centering\arraybackslash\hsize=3.0\hsize}X|}{\textbf{\laghu{नि}\laghu{क}\laghu{ट}} \par (3)} &
\multicolumn{2}{>{\centering\arraybackslash\hsize=2.0\hsize}X|}{\textbf{\laghu{न}\laghu{हिं}} \par (2)} &
\multicolumn{4}{>{\centering\arraybackslash\hsize=4.0\hsize}X|}{\textbf{\guru{आ}\guru{वै}} \par (4)} \\
\hline
\end{tabularx}
\vspace{-1.5em}
\end{table}

\begin{table}[h!]
\centering
\renewcommand{\arraystretch}{1.2}
\setlength{\tabcolsep}{0pt}
\begin{tabularx}{\textwidth}{|*{16}{>{\centering\arraybackslash\hsize=1.0\hsize}X|}}
\hline
\multicolumn{6}{|>{\centering\arraybackslash\hsize=6.0\hsize}X|}{\textbf{\laghu{म}\guru{हा}\guru{बी}\laghu{र}} \par (6)} &
\multicolumn{2}{>{\centering\arraybackslash\hsize=2.0\hsize}X|}{\textbf{\laghu{ज}\laghu{ब}} \par (2)} &
\multicolumn{3}{>{\centering\arraybackslash\hsize=3.0\hsize}X|}{\textbf{\guru{ना}\laghu{म}} \par (3)} &
\multicolumn{5}{>{\centering\arraybackslash\hsize=5.0\hsize}X|}{\textbf{\laghu{सु}\guru{ना}\guru{वै}} \par (5)} \\
\hline
\end{tabularx}
\vspace{-1.5em}
\end{table}

\vspace{0.5em}
\textbf{ \deva{सरल-संस्कृतम्}:}\\
 \deva{भूताः पिशाचाः च समीपम् अपि न आगच्छन्ति}\\
 \deva{यदा कोऽपि "महावीर" इति भवतः नाम श्रावयति (जपति)॥}\\

Ghosts and evil spirits do not even come near\\
whenever your name 'Mahaveer' is chanted aloud.


\vspace{1em}
\newpage
\textbf{\deva{प्रतिपदार्थः} (Meaning of each word):}
\begin{table}[h!]
\centering
\renewcommand{\arraystretch}{1.2}
\begin{tabularx}{\textwidth}{@{} l l X X @{}}
\toprule
\textbf{अवधी} & \textbf{संस्कृतम्} & \textbf{English} & \textbf{Grammar \& Notes} \\
\midrule
\deva{भूत} / \telu{భూత} / \textit{bhūta}  &  \deva{भूताः}  &  Ghosts  &   \\
\deva{पिसाच} / \telu{పిసాచ} / \textit{pisāca}  &  \deva{पिशाचाः}  &  Evil spirits / Ghouls  & \\ 
\deva{निकट} / \telu{నికట} / \textit{nikaṭa}  &  \deva{निकटम् / समीपम्}  &  Near / Close  &   \\
\deva{नहिं आवै} / \telu{నహిఁ ఆవై} / \textit{nahĩ āvai}  &  \deva{न आगच्छन्ति}  &  Do not come  &   \\
\deva{महाबीर} / \telu{మహాబీర} / \textit{mahābīra}  &  \deva{महावीरः}  &  Mahaveer (Great Hero)  & \\ 
\deva{जब} / \telu{జబ} / \textit{jaba}  &  \deva{यदा}  &  When  & \\ 
\deva{नाम} / \telu{నామ} / \textit{nāma}  &  \deva{नाम}  &  Name  &   \\
\deva{सुनावै} / \telu{సునావై} / \textit{sunāvai}  &  \deva{श्रावयति}  &  Is heard / recited  & Causative verb \\
\bottomrule
\end{tabularx}
\end{table}

\newpage
